\section{Definitions and Threat Model}

\begin{table*}[t]
\caption{Core terms.}
\label{tab:core-terms}
\centering
\scriptsize
\begin{tabularx}{\textwidth}{L{0.20\textwidth} Y}
\toprule
\JKTableHeader
\textbf{Term} & \textbf{Operational meaning} \\
\midrule
Vibe artifact & Repository residue that makes a change plausible but not auditable: ownerless paths, unmapped proof, hand-edited generated zones, stale contracts, false-green tests, missing security evidence, overbroad agent authority, unproven UI changes, stale waivers, or review claims without receipts. \\
Vibe-code drift & The measurable gap between claimed engineering intent and evidence attached to a change. \\
Continuous repository alignment & The practice of keeping ownership, edit scope, proof lanes, generated boundaries, permissions, evidence, repair queues, and conformance metadata synchronized with each merge. \\
Evidence debt & Missing or stale merge evidence that accumulates like technical debt: absent owner routes, unmapped proof lanes, unvalidated generated boundaries, undocumented waivers, or receipts that cannot prove the current change. \\
Review subjectivity gap & Reliance on taste, memory, seniority, or local convention where a rule ID, owner route, proof lane, receipt, or explicit waiver should exist. \\
Real-time repository conformance & A repeatable local/PR/CI loop in which each changed path is evaluated against current repository policy before merge; it is not a claim of an omniscient daemon. \\
Proof freshness & A receipt is current for the changed-path scope, commit, tool identity, command identity, and declared lane it claims to prove. \\
Receipt validity & A receipt is replayable or inspectable enough to bind command, cwd, inputs, tool version, result, artifacts, and digest to the current merge witness. \\
Policy authority & The ordered source of truth used to decide which instructions, maps, lanes, and adapter claims are trusted. \\
Adapter drift & Divergence between provider instruction files and canonical repository policy. \\
Evidence bundle & A versioned group of reports, receipts, logs, screenshots, ARIA snapshots, digests, and witness artifacts used to audit a claim. \\
Conformance claim & A machine-readable statement that a repository meets a named Jankurai level at a commit under a standard, auditor, and schema version. \\
Proof receipt & A durable record of a proof command, tool identity, inputs, result, and artifact path. \\
Merge witness & A versioned artifact binding changed paths, owner route, proof receipts, score delta, missing-evidence decision, artifact digests, commit identity, and tool version. \\
Generated zone & An output path repaired by changing the declared source contract or generator and regenerating, not by hand edit. \\
Repair queue & Ordered findings with rule IDs, owners, proof lanes, rerun commands, evidence requirements, and bounded next actions. \\
Reference profile & A non-normative stack or implementation profile that can satisfy the standard without defining it. \\
\bottomrule
\end{tabularx}
\end{table*}

\emph{Vibe coding} is any development mode where plausible code is accepted without explicit ownership, proof routing, generated-boundary checks, security gates, UX evidence, and repair evidence. Jankurai does not merely criticize vibe coding. It serializes vibe artifacts, assigns stable rule IDs, routes them to proof lanes, and reduces them through repair receipts.

\begin{table*}[t]
\caption{Threat model for auditable merge.}
\label{tab:threat-model}
\centering
\scriptsize
\begin{tabularx}{\textwidth}{L{0.18\textwidth} Y Y}
\toprule
\JKTableHeader
\textbf{Threat} & \textbf{Risk} & \textbf{Jankurai control} \\
\midrule
False plausibility & Agent or human patch looks idiomatic but violates architecture, data, security, or runtime expectations. & Stable rule IDs, owner maps, proof lanes, generated-zone policy, and merge witness. \\
Stale or copied claims & A stale artifact is attached to a commit, path scope, tool, or command that did not run. & Freshness checks, artifact digests, validity rules, and witness binding. \\
Adapter drift & Tool-specific instruction files widen permissions or contradict canonical policy. & Source-of-authority hierarchy, adapter verification, policy digests, and generated mirrors. \\
Flaky proof lanes & Intermittent tests hide real failures or make proof receipts unreliable. & Lane contracts with timeout, environment, retry policy, artifact paths, and failure classification. \\
Emergency bypass & Urgent work merges without owner, security, release, or proof evidence. & Expiring waivers, compensating controls, release-fail decisions, and repair receipts. \\
Evidence leakage & Screenshots, logs, traces, prompts, or artifacts expose secrets or customer data. & Redaction policy, artifact allowlists, secret scans, retention bounds, and privacy receipts. \\
CI downgrade & A workflow disables proof, widens permissions, removes artifacts, or changes branch gates. & CI hardening rule, workflow diff checks, pinned actions, and required evidence uploads. \\
Malicious MCP or tool schema & External tool output or schema attempts to change trusted policy or execute unsafe commands. & Untrusted-content isolation, schema validation, permission receipts, and no direct execution of generated text. \\
Generated test self-confirmation & Generated code and generated tests confirm the same wrong assumption. & Owner-owned acceptance criteria, negative fixtures, semantic assertions, and changed-path routing. \\
Subjective or rubber-stamp review & Review approval relies on taste, memory, or plausibility instead of reproducible evidence. & Stable rule IDs, human-review receipt rule, receipt requirement, explicit waiver or review receipt, and merge witness decision record. \\
Screenshots or logs leaking secrets & Browser and observability evidence captures tokens, PII, prompt text, or private data. & UX/security evidence redaction, hash-based references, and artifact privacy status. \\
Out of scope & Full semantic correctness, malicious compiler/runtime integrity, and organizational governance outside the repository. & Explicit limitations, signed-receipt future work, and no claim of total program proof. \\
\bottomrule
\end{tabularx}
\end{table*}

The standard does not claim every repository using AI must adopt Jankurai. It claims a narrower and testable thing: a repository claiming Jankurai conformance must expose enough machine-readable structure for the claim to be audited. Conformance is a public interface, not a vibe.
